\section{Conclusion}
\label{sec:conclusion}

This paper studies a two-stage wireless charging facility-location problem for electric-bus systems under service-plan uncertainty.
The model treats SWC and DWC deployment as strategic here-and-now decisions, while route-level battery selection and representative-day charging operations adapt after each service-plan scenario is realized.
The resulting recourse structure captures how wireless charging opportunities affect battery requirements and residual demand for limited shared plug-in fast chargers.

The main methodological contribution is an exact decomposition framework for this shared-resource mixed-integer recourse structure.
The method uses an exact lifted reformulation that keeps scenario-dependent battery choices adaptive, a nested scenario--route relaxation--verification decomposition, and joint-resource direction-selective monotonicity cuts.
Phase 1 route-level LP relaxations obtain valid lower bounds and cuts, Phase 2 restores exact mixed-integer operational verification, and the strengthened full method applies monotonicity cuts where they are valid for exact recourse.
Under exact master and recourse separation, the complete procedure converges to an optimal solution of the original two-stage formulation.

The computational study is designed to show how each algorithmic component contributes to scalability.
The base planning instance illustrates the type of wireless facility layout and adaptive battery decisions produced by the model.
The ablation, scenario-scaling, and capacity-coupling experiments then evaluate when the direct solver and simpler decomposition variants become insufficient and when the full method remains effective.

Several limitations define natural extensions.
The current experiments use constructed service-plan scenarios, treat plug-in fast-charging infrastructure as exogenous, and do not endogenize the feedback from battery weight to route-level energy consumption.
Future work can develop richer scenario-generation models, jointly optimize wireless and plug-in charging expansion, incorporate operational uncertainty, and apply the lifted two-phase decomposition idea to other stochastic infrastructure-planning problems with shared-resource recourse.
